\chapter{Polynomial Mealy Spans and Guarded Program Structure}
\label{ch:polynomial-mealy-spans-guarded-programs}

\section{Guarded Mealy structure as the reason for polynomial currying}
\label{sec:guarded-mealy-structure-polynomial-currying}

The previous chapter associated to every Mealy morphism
\[
(P,\delta,\omega):\mathbb A\rightsquigarrow\mathbb B
\]
a primitive execution span
\[
P
\xleftarrow{\pi_2}
A_1\times_{t_A,A_0,p}P
\xrightarrow{\delta}
P
\]
and hence an internal program category
\[
\mathsf{Prog}_\Phi.
\]
This span records individual executions
\[
(a,x),
\]
where \(a\) is an admissible input arrow at the state \(x\).

The span semantics already supplies the dynamic-logic structure:
tests, restricted programs, finite guarded choice, sequential composition,
path iteration, diamonds, and boxes. The purpose of the present chapter is
not to replace that span semantics. Rather, it is to expose the
\emph{guarded one-step response profile} hidden inside a Mealy state.

The point is that a Mealy state does not merely admit existential
one-step executions. It determines a total response profile:
for every admissible input query, it specifies both an output arrow and a
continuation state. Polynomial functors are the natural categorical
language for this curried profile. This use of polynomials follows the
standard shape-and-positions reading of polynomial functors and containers
\cite{GambinoKockPolynomialFunctorsMonads,AbbottAltenkirchGhaniContainers}.

Thus the guiding slogan of the chapter is:
\[
\text{dynamic logic is span-based,}
\qquad
\text{guarded Mealy response is polynomial.}
\]
Flattening the polynomial coalgebra recovers the execution span. Pulling
guarded polynomial predicate liftings back along the coalgebra recovers
the guarded span modalities.

The guarded dynamic logic developed below is not a second, independent
semantics. It is the span-based dynamic logic of
Chapter~\ref{ch:mealy-program-categories-dynamic-logic} applied to the
flattened execution span of the polynomial Mealy coalgebra, with state,
input, output, and interface predicates used as guards.

Let
\[
\mathbb A=(A_0,A_1,s_A,t_A,e_A,m_A)
\]
and
\[
\mathbb B=(B_0,B_1,s_B,t_B,e_B,m_B)
\]
be internal categories in a category \(\mathcal E\).

For the polynomial construction, assume that the products
\[
A_0\times B_0,
\qquad
A_1\times B_1
\]
exist, and that the dependent products used below exist. A sufficient
global hypothesis is that \(\mathcal E\) is locally cartesian closed with
finite products. For guarded diamond semantics we additionally assume
regularity. For finite guarded choice we assume coherence. For boxes and
guarded iteration we use the corresponding first-order Heyting and
\(\omega\)-iterative hypotheses from
Chapter~\ref{ch:mealy-program-categories-dynamic-logic}.

Let
\[
S:=A_0\times B_0.
\]
A Mealy carrier
\[
A_0\xleftarrow{p}P\xrightarrow{q}B_0
\]
is equivalently an object over \(S\):
\[
(p,q):P\to S.
\]
This is displayed by the commuting diagram
\[
\begin{tikzcd}[column sep=large,row sep=large]
&
P
  \arrow[dl, "p"']
  \arrow[dr, "q"]
  \arrow[d, "{(p,q)}"]
&
\\
A_0
&
S=A_0\times B_0
  \arrow[l, "\pi_A"]
  \arrow[r, "\pi_B"']
&
B_0 .
\end{tikzcd}
\]
Thus a state \(x\in P\) lies over an interface pair
\[
(p(x),q(x))\in A_0\times B_0.
\]

The primitive execution object is the pullback
\[
\begin{tikzcd}[column sep=large,row sep=large]
M_1=A_1\times_{t_A,A_0,p}P
  \arrow[r, "\pi_1"]
  \arrow[d, "\pi_2"']
&
A_1
  \arrow[d, "t_A"]
\\
P
  \arrow[r, "p"']
&
A_0 .
\end{tikzcd}
\]
The primitive execution span is
\[
\begin{tikzcd}[column sep=large]
P
&
M_1
  \arrow[l, "\pi_2"']
  \arrow[r, "\delta"]
&
P .
\end{tikzcd}
\]

For a state
\[
x\in P_{(v,y)}
\]
and an admissible input arrow
\[
a:u\to v
\]
of \(\mathbb A\), the Mealy structure produces
\[
\delta(a,x)\in P
\]
and
\[
\omega(a,x)\in B_1.
\]
The typing equations say
\[
p(\delta(a,x))=u,
\]
\[
s_B(\omega(a,x))=q(\delta(a,x)),
\]
and
\[
t_B(\omega(a,x))=q(x)=y.
\]
Thus the one-step response to \(a\) is a pair
\[
\bigl(\omega(a,x),\delta(a,x)\bigr)
\in
\sum_{\beta:y'\to y}P_{(u,y')}.
\]

The typing is summarized by
\[
\begin{tikzcd}[column sep=large,row sep=large]
q(\delta(a,x))
  \arrow[r, "{\omega(a,x)}"]
&
q(x)
\\
s_A(a)=u
  \arrow[r, "a"']
&
v=p(x) .
\end{tikzcd}
\]

Consequently, the whole one-step behaviour of \(x\) is a total response
profile
\[
a
\longmapsto
\bigl(\omega(a,x),\delta(a,x)\bigr)
\]
of type
\[
\prod_{a:u\to v}
\sum_{\beta:y'\to y}
P_{(u,y')}.
\]

The central construction of this chapter is the canonical polynomial
functor whose fibre over \((v,y)\) is
\[
\prod_{a:u\to v}
\sum_{\beta:y'\to y}
X_{(u,y')}.
\]
A Mealy morphism is then precisely a coalgebra for this functor satisfying
the identity and multiplicativity equations induced by the internal
category structures of \(\mathbb A\) and \(\mathbb B\).

\section{Polynomial spans as dependent response transformers}
\label{sec:polynomial-spans-dependent-response-transformers}

We recall the polynomial notation needed for the construction.

For a morphism
\[
f:X\to Y
\]
write
\[
f^\ast:\mathcal E/Y\to\mathcal E/X
\]
for pullback along \(f\),
\[
\Sigma_f:\mathcal E/X\to\mathcal E/Y
\]
for postcomposition with \(f\), and
\[
\Pi_f:\mathcal E/X\to\mathcal E/Y
\]
for the right adjoint to \(f^\ast\), when it exists:
\[
\Sigma_f\dashv f^\ast\dashv \Pi_f.
\]
Equivalently, one has the adjunction string
\[
\begin{tikzcd}[column sep=large]
\mathcal E/X
  \arrow[r, shift left=1.1ex, "\Sigma_f"]
  \arrow[r, shift right=1.1ex, swap, "\Pi_f"]
&
\mathcal E/Y
  \arrow[l, "f^\ast" description] .
\end{tikzcd}
\]
We use the standard slice-categorical notation for dependent sums,
pullback, and dependent products; see, for example,
\cite{JacobsCategoricalLogicTypeTheory,GambinoKockPolynomialFunctorsMonads}.

\begin{definition}[Polynomial span]
\label{def:polynomial-span}
A \textbf{polynomial span} from \(I\) to \(J\) is a diagram
\[
\begin{tikzcd}[column sep=large]
I
&
E
  \arrow[l, "s"']
  \arrow[r, "p"]
&
B
  \arrow[r, "t"]
&
J .
\end{tikzcd}
\]
When the required dependent product exists, it induces a functor
\[
\Sigma_t\Pi_p s^\ast:
\mathcal E/I\to\mathcal E/J.
\]
\end{definition}

This is the usual polynomial-functor construction over a locally
cartesian closed base, written in slice-functor form
\[
\Sigma_t\Pi_p s^\ast
\]
\cite{GambinoKockPolynomialFunctorsMonads}.

Thus a polynomial span is read as the composite of slice functors
\[
\begin{tikzcd}[column sep=large]
\mathcal E/I
  \arrow[r, "s^\ast"]
&
\mathcal E/E
  \arrow[r, "\Pi_p"]
&
\mathcal E/B
  \arrow[r, "\Sigma_t"]
&
\mathcal E/J .
\end{tikzcd}
\]

In generalized-element notation, if
\[
X\to I
\]
is an object over \(I\), then
\[
\left(\Sigma_t\Pi_p s^\ast X\right)_j
\cong
\sum_{b\in B_j}
\prod_{e\in E_b}
X_{s(e)}.
\]

\begin{remark}[Ordinary spans as the linear case]
\label{rem:ordinary-spans-linear-polynomial-case}
An ordinary span
\[
I\xleftarrow{s}E\xrightarrow{t}J
\]
is the special case of a polynomial span in which the object of shapes is
the object of positions itself:
\[
\begin{tikzcd}[column sep=large]
I
&
E
  \arrow[l, "s"']
  \arrow[r, "1_E"]
&
E
  \arrow[r, "t"]
&
J .
\end{tikzcd}
\]
The induced functor is
\[
\Sigma_t\Pi_{1_E}s^\ast
\cong
\Sigma_t s^\ast.
\]
Thus ordinary spans are the linear case: each execution witness has
exactly one continuation position. The polynomial structure below arises
by replacing this single continuation position by a dependent family of
query positions.
\end{remark}

\begin{remark}[Linear polynomial orientation for execution spans]
\label{rem:linear-polynomial-execution-orientation}
For an execution span
\[
P\xleftarrow{\ell}R\xrightarrow{r}P
\]
with \(\ell\) the current-state/source leg and \(r\) the
continuation-state/target leg, the corresponding linear polynomial
response transformer uses
\[
s=r,
\qquad
t=\ell,
\]
because \(r\) records the continuation state and \(\ell\) records the
current state. This is the same orientation used below for the Mealy
polynomial: the polynomial left leg records continuation interface, while
the final right leg records current interface.
\end{remark}

\section{Queries, responses, and the canonical Mealy polynomial}
\label{sec:queries-responses-canonical-mealy-polynomial}

We now construct the polynomial span forced by the total guarded Mealy
response type.

Let
\[
S=A_0\times B_0
\]
with projections
\[
\pi_A:S\to A_0,
\qquad
\pi_B:S\to B_0.
\]

\subsection{Queries}

Define the object of admissible input queries by
\[
Q_{\mathbb A}:=
A_1\times_{t_A,A_0,\pi_A}S.
\]
It is the pullback
\[
\begin{tikzcd}[column sep=large,row sep=large]
Q_{\mathbb A}
  \arrow[r, "\overline a"]
  \arrow[d, "\kappa"']
&
A_1
  \arrow[d, "t_A"]
\\
S
  \arrow[r, "\pi_A"']
&
A_0 .
\end{tikzcd}
\]
Here \(\overline a\) denotes the projection to \(A_1\), and
\[
\kappa:Q_{\mathbb A}\to S
\]
records the current interface.

A generalized element of \(Q_{\mathbb A}\) is written
\[
(a,y),
\]
where
\[
a:u\to v
\]
is an arrow of \(\mathbb A\), and \(y\in B_0\). The current interface is
\[
(v,y).
\]
Thus
\[
\kappa(a,y)=(t_A(a),y).
\]
There is also a map recording the next input interface:
\[
\lambda_A:Q_{\mathbb A}\to A_0,
\qquad
\lambda_A(a,y)=s_A(a).
\]
Formally,
\[
\lambda_A:=s_A\overline a.
\]

Thus
\[
\kappa:Q_{\mathbb A}\to S
\]
is the object of input queries available at each current interface.

\subsection{Responses}

A response to a query
\[
(a:u\to v,\ y)
\]
is an output arrow of \(\mathbb B\)
\[
\beta:y'\to y.
\]
The next output interface is then
\[
y'=s_B(\beta).
\]

Define
\[
O_{\mathbb A,\mathbb B}:=A_1\times B_1.
\]
A generalized element is written
\[
(a,\beta).
\]

Define
\[
r:O_{\mathbb A,\mathbb B}\to Q_{\mathbb A}
\]
to be the unique map satisfying
\[
\overline a\,r=\pi_{A_1},
\qquad
\kappa r=(t_A\pi_{A_1},t_B\pi_{B_1}).
\]
In generalized-element notation,
\[
r(a,\beta)=(a,t_B(\beta)).
\]
This is well-defined because the associated current interface is
\[
(t_A(a),t_B(\beta)).
\]

There is also a map
\[
n:O_{\mathbb A,\mathbb B}\to S
\]
given by
\[
n(a,\beta)=(s_A(a),s_B(\beta)).
\]
Equivalently,
\[
n=(s_A\pi_{A_1},s_B\pi_{B_1}).
\]
This records the continuation interface after consuming \(a\) and emitting
\(\beta\).

These maps are summarized by
\[
\begin{tikzcd}[column sep=large,row sep=large]
O_{\mathbb A,\mathbb B}=A_1\times B_1
  \arrow[r, "r"]
  \arrow[d, "n"']
&
Q_{\mathbb A}
  \arrow[d, "\kappa"]
\\
S
&
S .
\end{tikzcd}
\]
The square is not asserted to be a pullback. It records two interface
maps: \(r\) remembers the query determined by the current output endpoint,
and \(n\) remembers the continuation interface.

Thus we have a diagram
\[
\begin{tikzcd}[column sep=large]
S
&
O_{\mathbb A,\mathbb B}
  \arrow[l, "n"']
  \arrow[r, "r"]
&
Q_{\mathbb A}
  \arrow[r, "\kappa"]
&
S .
\end{tikzcd}
\]

\subsection{Total output policies}

Define the object of total output policies by the dependent product
\[
K_{\mathbb A,\mathbb B}
:=
\Pi_\kappa
\bigl(r:O_{\mathbb A,\mathbb B}\to Q_{\mathbb A}\bigr).
\]
This is an object over \(S\). We write
\[
t_{\mathbb A,\mathbb B}:K_{\mathbb A,\mathbb B}\to S
\]
for its structure map.

The construction is summarized by
\[
\begin{tikzcd}[column sep=large,row sep=large]
O_{\mathbb A,\mathbb B}
  \arrow[d, "r"']
&
K_{\mathbb A,\mathbb B}
  \arrow[d, "t_{\mathbb A,\mathbb B}"]
\\
Q_{\mathbb A}
  \arrow[r, "\kappa"']
&
S ,
\end{tikzcd}
\]
where \(K_{\mathbb A,\mathbb B}\) is the dependent product of \(r\) along
\(\kappa\).

A generalized element of \(K_{\mathbb A,\mathbb B}\) over
\((v,y)\in S\) is an assignment
\[
\sigma:
\{a:u\to v\}
\longrightarrow
\{\beta:y'\to y\}.
\]
Thus \(\sigma\) assigns to every admissible input arrow \(a:u\to v\) an
output arrow with target \(y\).

Now form the pullback
\[
E_{\mathbb A,\mathbb B}
:=
Q_{\mathbb A}\times_{S}K_{\mathbb A,\mathbb B},
\]
where the maps to \(S\) are
\[
\kappa:Q_{\mathbb A}\to S
\]
and
\[
t_{\mathbb A,\mathbb B}:K_{\mathbb A,\mathbb B}\to S.
\]
Thus
\[
\begin{tikzcd}[column sep=large,row sep=large]
E_{\mathbb A,\mathbb B}
  \arrow[r, "\epsilon_Q"]
  \arrow[d, "p_{\mathbb A,\mathbb B}"']
&
Q_{\mathbb A}
  \arrow[d, "\kappa"]
\\
K_{\mathbb A,\mathbb B}
  \arrow[r, "t_{\mathbb A,\mathbb B}"']
&
S
\end{tikzcd}
\]
is a pullback square. The map
\[
p_{\mathbb A,\mathbb B}:E_{\mathbb A,\mathbb B}\to K_{\mathbb A,\mathbb B}
\]
is the second projection.

The counit of
\[
\kappa^\ast\dashv\Pi_\kappa
\]
gives an evaluation map
\[
\operatorname{ev}:E_{\mathbb A,\mathbb B}\to O_{\mathbb A,\mathbb B}
\]
over \(Q_{\mathbb A}\):
\[
\begin{tikzcd}[column sep=large,row sep=large]
E_{\mathbb A,\mathbb B}
  \arrow[r, "\operatorname{ev}"]
  \arrow[d, "\epsilon_Q"']
&
O_{\mathbb A,\mathbb B}
  \arrow[d, "r"]
\\
Q_{\mathbb A}
  \arrow[r, equal]
&
Q_{\mathbb A}.
\end{tikzcd}
\]
In generalized-element notation,
\[
\operatorname{ev}(a,\sigma)=(a,\sigma(a)).
\]

Composing with
\[
n:O_{\mathbb A,\mathbb B}\to S
\]
gives
\[
s_{\mathbb A,\mathbb B}
:=
n\circ\operatorname{ev}
:
E_{\mathbb A,\mathbb B}\to S.
\]
Thus \(s_{\mathbb A,\mathbb B}\) records the continuation interface.

We have therefore constructed a polynomial span over \(S\):
\[
\begin{tikzcd}[column sep=large]
S
&
E_{\mathbb A,\mathbb B}
  \arrow[l, "s_{\mathbb A,\mathbb B}"']
  \arrow[r, "p_{\mathbb A,\mathbb B}"]
&
K_{\mathbb A,\mathbb B}
  \arrow[r, "t_{\mathbb A,\mathbb B}"]
&
S .
\end{tikzcd}
\]

\begin{definition}[Mealy polynomial span]
\label{def:mealy-polynomial-span}
The polynomial span
\[
\begin{tikzcd}[column sep=large]
S
&
E_{\mathbb A,\mathbb B}
  \arrow[l, "s_{\mathbb A,\mathbb B}"']
  \arrow[r, "p_{\mathbb A,\mathbb B}"]
&
K_{\mathbb A,\mathbb B}
  \arrow[r, "t_{\mathbb A,\mathbb B}"]
&
S ,
\end{tikzcd}
\qquad
S=A_0\times B_0,
\]
is called the \textbf{Mealy polynomial span} associated to the pair
\[
(\mathbb A,\mathbb B).
\]
\end{definition}

\begin{remark}[Polynomial orientation]
\label{rem:polynomial-orientation-mealy}
The symbols
\[
s_{\mathbb A,\mathbb B}
\qquad\text{and}\qquad
t_{\mathbb A,\mathbb B}
\]
belong to the polynomial orientation, not to the execution-span
orientation. Here
\[
t_{\mathbb A,\mathbb B}
\]
records the current interface of a total response policy, while
\[
s_{\mathbb A,\mathbb B}
\]
records the continuation interface of a queried position.
\end{remark}

\section{The Mealy polynomial functor}
\label{sec:mealy-polynomial-functor}

The Mealy polynomial span induces an endofunctor
\[
\mathcal M_{\mathbb A,\mathbb B}:\mathcal E/S\to\mathcal E/S
\]
defined by
\[
\mathcal M_{\mathbb A,\mathbb B}(X)
:=
\Sigma_{t_{\mathbb A,\mathbb B}}
\Pi_{p_{\mathbb A,\mathbb B}}
s_{\mathbb A,\mathbb B}^{\ast}X.
\]
That is,
\[
\begin{tikzcd}[column sep=large]
\mathcal E/S
  \arrow[r, "s_{\mathbb A,\mathbb B}^{\ast}"]
&
\mathcal E/E_{\mathbb A,\mathbb B}
  \arrow[r, "\Pi_{p_{\mathbb A,\mathbb B}}"]
&
\mathcal E/K_{\mathbb A,\mathbb B}
  \arrow[r, "\Sigma_{t_{\mathbb A,\mathbb B}}"]
&
\mathcal E/S .
\end{tikzcd}
\]

\begin{proposition}[Fibre calculation for the Mealy polynomial]
\label{prop:mealy-polynomial-fibre}
Let
\[
X\to S=A_0\times B_0
\]
be an object over \(S\). For every generalized interface pair \((v,y)\),
there is a canonical fibrewise description
\[
\mathcal M_{\mathbb A,\mathbb B}(X)_{(v,y)}
\cong
\prod_{a:u\to v}
\sum_{\beta:y'\to y}
X_{(u,y')}.
\]
\end{proposition}

\begin{proof}
A point of
\[
K_{\mathbb A,\mathbb B}
\]
over \((v,y)\) is a total output policy
\[
\sigma:
a:u\to v
\longmapsto
\sigma(a):y_a\to y.
\]
The object
\[
E_{\mathbb A,\mathbb B}
=
Q_{\mathbb A}\times_S K_{\mathbb A,\mathbb B}
\]
has, over such a policy \(\sigma\), one position for each input arrow
\[
a:u\to v.
\]
The map
\[
s_{\mathbb A,\mathbb B}:E_{\mathbb A,\mathbb B}\to S
\]
sends this position to the next interface
\[
(u,y_a)
=
(s_A(a),s_B(\sigma(a))).
\]
Therefore
\[
\Pi_{p_{\mathbb A,\mathbb B}}
s_{\mathbb A,\mathbb B}^{\ast}X
\]
assigns to \(\sigma\) the family of choices
\[
x_a\in X_{(u,y_a)}
\]
for every admissible input arrow \(a:u\to v\). Finally
\[
\Sigma_{t_{\mathbb A,\mathbb B}}
\]
regards the resulting data as lying over the current interface \((v,y)\).

Thus
\[
\mathcal M_{\mathbb A,\mathbb B}(X)_{(v,y)}
\cong
\sum_{\sigma}
\prod_{a:u\to v}
X_{(s_A(a),s_B(\sigma(a)))}.
\]
Choosing such a policy \(\sigma\) together with all the \(x_a\)'s is
equivalently, for every \(a:u\to v\), choosing a pair
\[
\bigl(\beta:y'\to y,\ x_a\in X_{(u,y')}\bigr).
\]
This uses the standard dependent distributivity isomorphism available in
a locally cartesian closed category:
\[
\sum_{\sigma:\prod_a B_a}
\prod_a X_{a,\sigma(a)}
\cong
\prod_a
\sum_{\beta\in B_a}X_{a,\beta}.
\]
Here \(B_a\) is the fibre of admissible output arrows
\[
\beta:y'\to y
\]
for the input query \(a:u\to v\). Hence
\[
\mathcal M_{\mathbb A,\mathbb B}(X)_{(v,y)}
\cong
\prod_{a:u\to v}
\sum_{\beta:y'\to y}
X_{(u,y')}.
\]
\end{proof}

\begin{remark}[Container reading of the fibre calculation]
\label{rem:container-reading-fibre-calculation}
The fibre calculation is the standard container-style reading of a
polynomial: an element consists of a shape, together with a dependent
choice of data at every position of that shape
\cite{AbbottAltenkirchGhaniContainers,GambinoKockPolynomialFunctorsMonads}.
\end{remark}

\begin{remark}[Local hypotheses]
\label{rem:local-polynomial-hypotheses}
If one does not assume that \(\mathcal E\) is globally locally cartesian
closed, then the preceding construction should be read under the local
assumption that the displayed dependent products and the dependent
distributivity isomorphisms used in the fibre calculation exist.
\end{remark}

\section{Polynomial currying of one-step Mealy data}
\label{sec:polynomial-currying-one-step-mealy-data}

Let
\[
P\to S=A_0\times B_0
\]
be an object over \(S\). Write the structure map as
\[
(p,q):P\to A_0\times B_0.
\]
Thus
\[
\begin{tikzcd}[column sep=large,row sep=large]
P
  \arrow[r, "{(p,q)}"]
  \arrow[dr, "p"']
  \arrow[drr, bend left=15, "q"]
&
S=A_0\times B_0
  \arrow[d, "\pi_A"]
  \arrow[dr, "\pi_B"]
&
\\
&
A_0
&
B_0 .
\end{tikzcd}
\]

Suppose first that we are given one-step Mealy data
\[
\delta:
A_1\times_{t_A,A_0,p}P\to P,
\]
\[
\omega:
A_1\times_{t_A,A_0,p}P\to B_1
\]
satisfying the typing equations
\begin{align}
p\circ\delta&=s_A\circ\pi_1,
\label{eq:poly-curry-pdelta}
\\
s_B\circ\omega&=q\circ\delta,
\label{eq:poly-curry-source-output}
\\
t_B\circ\omega&=q\circ\pi_2.
\label{eq:poly-curry-target-output}
\end{align}

The typing data are summarized by
\[
\begin{tikzcd}[column sep=large,row sep=large]
M_1=A_1\times_{t_A,A_0,p}P
  \arrow[r, "\delta"]
  \arrow[d, "\pi_1"']
  \arrow[dr, phantom, "\scriptstyle p\delta=s_A\pi_1"]
&
P
  \arrow[d, "p"]
\\
A_1
  \arrow[r, "s_A"']
&
A_0
\end{tikzcd}
\]
and
\[
\begin{tikzcd}[column sep=large,row sep=large]
M_1
  \arrow[r, "\omega"]
  \arrow[d, "\delta"']
&
B_1
  \arrow[d, "s_B"]
\\
P
  \arrow[r, "q"']
&
B_0
\end{tikzcd}
\qquad
\begin{tikzcd}[column sep=large,row sep=large]
M_1
  \arrow[r, "\omega"]
  \arrow[d, "\pi_2"']
&
B_1
  \arrow[d, "t_B"]
\\
P
  \arrow[r, "q"']
&
B_0 .
\end{tikzcd}
\]

For
\[
x\in P_{(v,y)}
\]
and
\[
a:u\to v,
\]
the pair
\[
\bigl(\omega(a,x),\delta(a,x)\bigr)
\]
lies in
\[
\sum_{\beta:y'\to y}P_{(u,y')}.
\]
Thus
\[
a
\longmapsto
\bigl(\omega(a,x),\delta(a,x)\bigr)
\]
is an element of
\[
\prod_{a:u\to v}
\sum_{\beta:y'\to y}
P_{(u,y')}.
\]
By Proposition~\ref{prop:mealy-polynomial-fibre}, this is an element of
\[
\mathcal M_{\mathbb A,\mathbb B}(P)_{(v,y)}.
\]

Hence the one-step Mealy data determine a morphism over \(S\)
\[
\widehat{\Phi}:P\to\mathcal M_{\mathbb A,\mathbb B}(P)
\]
defined by
\[
\widehat{\Phi}(x)(a)
=
\bigl(\omega(a,x),\delta(a,x)\bigr).
\]
The fact that this is a map over \(S\) is expressed by
\[
\begin{tikzcd}[column sep=large,row sep=large]
P
  \arrow[r, "\widehat{\Phi}"]
  \arrow[dr, "{(p,q)}"']
&
\mathcal M_{\mathbb A,\mathbb B}(P)
  \arrow[d]
\\
&
S .
\end{tikzcd}
\]

\begin{theorem}[Polynomial currying of one-step Mealy data]
\label{thm:polynomial-currying-one-step}
Let
\[
P\to A_0\times B_0
\]
be an object over the interface object. To give maps
\[
\delta:
A_1\times_{t_A,A_0,p}P\to P,
\qquad
\omega:
A_1\times_{t_A,A_0,p}P\to B_1
\]
satisfying the typing equations
\[
p\delta=s_A\pi_1,
\qquad
s_B\omega=q\delta,
\qquad
t_B\omega=q\pi_2
\]
is equivalently to give a coalgebra map
\[
c:P\to\mathcal M_{\mathbb A,\mathbb B}(P)
\]
in the slice \(\mathcal E/S\).
\end{theorem}

\begin{proof}
Given \((\delta,\omega)\), define
\[
c=\widehat{\Phi}
\]
by
\[
c(x)(a)=\bigl(\omega(a,x),\delta(a,x)\bigr).
\]
The typing equations are exactly the assertion that
\[
\omega(a,x):q(\delta(a,x))\to q(x)
\]
and
\[
\delta(a,x)\in P_{(s_A(a),q(\delta(a,x)))}.
\]
Thus the displayed pair belongs to the fibre
\[
\sum_{\beta:y'\to q(x)}
P_{(s_A(a),y')}.
\]
Therefore \(c\) is a morphism over \(S\).

Conversely, suppose
\[
c:P\to\mathcal M_{\mathbb A,\mathbb B}(P)
\]
is a coalgebra over \(S\). For
\[
x\in P_{(v,y)}
\]
and
\[
a:u\to v,
\]
write
\[
c(x)(a)=(\beta_a,x_a)
\]
with
\[
\beta_a:y_a\to y
\]
and
\[
x_a\in P_{(u,y_a)}.
\]
Define
\[
\omega(a,x):=\beta_a,
\qquad
\delta(a,x):=x_a.
\]
Then
\[
p(\delta(a,x))=u=s_A(a),
\]
\[
s_B(\omega(a,x))=y_a=q(\delta(a,x)),
\]
and
\[
t_B(\omega(a,x))=y=q(x).
\]
Hence the Mealy typing equations hold.

The two constructions are inverse by inspection.
\end{proof}

\section{Multiplicative Mealy coalgebras}
\label{sec:multiplicative-mealy-coalgebras}

The coalgebra map
\[
c:P\to\mathcal M_{\mathbb A,\mathbb B}(P)
\]
records the typed one-step response profile. The remaining Mealy axioms
say that this response profile respects identities and composition in the
input and output internal categories.

Write
\[
c(x)(a)=(\beta_a,x_a).
\]
Thus
\[
\beta_a:q(x_a)\to q(x),
\qquad
x_a\in P_{(s_A(a),q(x_a))}.
\]

\begin{definition}[Multiplicative Mealy coalgebra]
\label{def:multiplicative-mealy-coalgebra}
A coalgebra
\[
c:P\to\mathcal M_{\mathbb A,\mathbb B}(P)
\]
is called a \textbf{multiplicative Mealy coalgebra} if it satisfies the
following equations.

\begin{enumerate}
\item For every state \(x\in P\),
\[
c(x)(e_A(px))=(e_B(qx),x).
\]

\item For every composable pair
\[
a:u\to v,
\qquad
b:v\to w,
\]
and every
\[
x\in P_{(w,y)},
\]
if
\[
c(x)(b)=(\beta,x')
\]
and
\[
c(x')(a)=(\alpha,x'')
\]
then
\[
c(x)(b\circ a)=(\beta\circ\alpha,x'').
\]
Equivalently, since
\[
m_A(a,b)=b\circ a,
\qquad
m_B(\alpha,\beta)=\beta\circ\alpha,
\]
this condition is
\[
c(x)(m_A(a,b))
=
\bigl(m_B(\alpha,\beta),x''\bigr).
\]
\end{enumerate}
\end{definition}

Thus a multiplicative Mealy coalgebra is not merely a coalgebra for the
polynomial endofunctor \(\mathcal M_{\mathbb A,\mathbb B}\). It is a
coalgebra whose response profile is compatible with the identity and
composition operations of the internal categories \(\mathbb A\) and
\(\mathbb B\).

The composite response can be pictured as
\[
\begin{tikzcd}[column sep=large,row sep=large]
x
  \arrow[r, "{b/\beta}"]
&
x'
  \arrow[r, "{a/\alpha}"]
&
x''
\\
q(x)
&
q(x')
  \arrow[l, "\beta"']
&
q(x'')
  \arrow[l, "\alpha"']
\\
w=p(x)
&
v=p(x')
  \arrow[l, "b"']
&
u=p(x'')
  \arrow[l, "a"']
\end{tikzcd}
\]
so that the folded output is
\[
\beta\circ\alpha:q(x'')\to q(x).
\]

\begin{remark}[Formal reading of multiplicativity]
\label{rem:formal-reading-multiplicativity}
The multiplication equation is formally an equality of morphisms out of
the pullback object
\[
D
:=
\bigl(A_1\times_{t_A,A_0,s_A}A_1\bigr)
\times_{t_A\pi_2,A_0,p}
P.
\]
Equivalently,
\[
\begin{tikzcd}[column sep=large,row sep=large]
D
  \arrow[r]
  \arrow[d]
&
P
  \arrow[d, "p"]
\\
A_1\times_{t_A,A_0,s_A}A_1
  \arrow[r, "t_A\pi_2"']
&
A_0 .
\end{tikzcd}
\]
A generalized element of \(D\) is a triple
\[
(a,b,x)
\]
satisfying
\[
t_A(a)=s_A(b),
\qquad
t_A(b)=p(x).
\]
On this object, the two sides of the multiplicativity equation are the
one-step response to
\[
m_A(a,b)
\]
and the response obtained by first responding to \(b\) at \(x\), then
responding to \(a\) at the resulting state.
\end{remark}

\begin{theorem}[Mealy morphisms as multiplicative coalgebras]
\label{thm:mealy-morphisms-multiplicative-coalgebras}
Let
\[
P\to A_0\times B_0
\]
be an object over the interface object. To give a Mealy morphism
\[
(P,\delta,\omega):\mathbb A\rightsquigarrow\mathbb B
\]
with carrier \(P\) is equivalently to give a multiplicative Mealy
coalgebra
\[
c:P\to\mathcal M_{\mathbb A,\mathbb B}(P).
\]
\end{theorem}

\begin{proof}
By Theorem~\ref{thm:polynomial-currying-one-step}, the typed one-step data
\[
(\delta,\omega)
\]
are equivalent to a coalgebra
\[
c:P\to\mathcal M_{\mathbb A,\mathbb B}(P)
\]
with
\[
c(x)(a)=(\omega(a,x),\delta(a,x)).
\]

The Mealy unit equations
\[
\delta(e_A(px),x)=x,
\qquad
\omega(e_A(px),x)=e_B(qx)
\]
are exactly
\[
c(x)(e_A(px))=(e_B(qx),x).
\]

For composition, let
\[
a:u\to v,
\qquad
b:v\to w,
\qquad
p(x)=w.
\]
If
\[
c(x)(b)=(\beta,x')
\]
and
\[
c(x')(a)=(\alpha,x'')
\]
then
\[
\beta=\omega(b,x),
\qquad
x'=\delta(b,x),
\]
and
\[
\alpha=\omega(a,\delta(b,x)),
\qquad
x''=\delta(a,\delta(b,x)).
\]
The coalgebra multiplication equation
\[
c(x)(b\circ a)=(\beta\circ\alpha,x'')
\]
is therefore exactly the pair of Mealy multiplication equations
\[
\delta(b\circ a,x)
=
\delta(a,\delta(b,x))
\]
and
\[
\omega(b\circ a,x)
=
\omega(b,x)\circ\omega(a,\delta(b,x)).
\]
Hence the Mealy axioms are precisely the multiplicativity equations for
\(c\).
\end{proof}

\section{Flattening the coalgebra back to the execution category}
\label{sec:flattening-coalgebra-execution-category}

Let
\[
c:P\to\mathcal M_{\mathbb A,\mathbb B}(P)
\]
be a multiplicative Mealy coalgebra. By
Theorem~\ref{thm:polynomial-currying-one-step}, it determines maps
\[
\delta:
A_1\times_{t_A,A_0,p}P\to P,
\qquad
\omega:
A_1\times_{t_A,A_0,p}P\to B_1.
\]
Hence it determines the primitive execution span
\[
\begin{tikzcd}[column sep=large]
P
&
A_1\times_{t_A,A_0,p}P
  \arrow[l, "\pi_2"']
  \arrow[r, "\delta"]
&
P .
\end{tikzcd}
\]

Equivalently, the flattened execution object is obtained as the pullback
of the position object of the polynomial response profile along the
coalgebra:
\[
\begin{tikzcd}[column sep=large,row sep=large]
A_1\times_{t_A,A_0,p}P
  \arrow[r]
  \arrow[d, "\pi_2"']
&
\mathrm{Pos}_P
  \arrow[d, "\rho"]
\\
P
  \arrow[r, "c"']
&
\mathcal M_{\mathbb A,\mathbb B}(P).
\end{tikzcd}
\]
Under this identification, the evaluations
\[
\operatorname{in},
\qquad
\operatorname{out},
\qquad
\operatorname{next}
\]
become respectively
\[
\pi_1,
\qquad
\omega,
\qquad
\delta.
\]

The multiplicativity equations of \(c\) are exactly the equations used in
Chapter~\ref{ch:mealy-program-categories-dynamic-logic} to make this span
into a monoid object in
\[
\operatorname{EndSpan}_{\mathcal E}(P).
\]

\begin{theorem}[Program category recovered from the polynomial coalgebra]
\label{thm:program-category-recovered-polynomial-coalgebra}
Let
\[
c:P\to\mathcal M_{\mathbb A,\mathbb B}(P)
\]
be the multiplicative Mealy coalgebra corresponding to a Mealy morphism
\[
(P,\delta,\omega):\mathbb A\rightsquigarrow\mathbb B.
\]
Then the primitive execution span recovered from \(c\) is exactly
\[
\mathbf m_\Phi
=
\left(
P
\xleftarrow{\pi_2}
A_1\times_{t_A,A_0,p}P
\xrightarrow{\delta}
P
\right),
\]
and the internal category it generates is the program category
\[
\mathsf{Prog}_\Phi
\]
of Chapter~\ref{ch:mealy-program-categories-dynamic-logic}.
\end{theorem}

\begin{proof}
The coalgebra \(c\) recovers \(\delta\) by
\[
c(x)(a)=(\omega(a,x),\delta(a,x)).
\]
Therefore the recovered primitive execution span has arrow object
\[
A_1\times_{t_A,A_0,p}P
\]
with source \(\pi_2\) and target \(\delta\), exactly as in the definition
of \(\mathbf m_\Phi\).

The identity equation for \(c\) gives
\[
\delta(e_A(px),x)=x,
\]
so the unit map
\[
e_M(x)=(e_A(px),x)
\]
is a map of spans
\[
1_P\Rightarrow \mathbf m_\Phi.
\]
Indeed,
\[
\pi_2 e_M=1_P,
\qquad
\delta e_M=1_P.
\]

The multiplication equation for \(c\) gives
\[
\delta(b\circ a,x)=\delta(a,\delta(b,x)),
\]
so the multiplication
\[
m_M\bigl((b,x),(a,\delta(b,x))\bigr)
=
(b\circ a,x)
\]
is a map of spans
\[
\mathbf m_\Phi;\mathbf m_\Phi\Rightarrow\mathbf m_\Phi.
\]
The relevant apex map is
\[
m_M:
M_1\times_{t_M,P,s_M}M_1\to M_1,
\]
and the source and target preservation conditions are
\[
s_M m_M=s_M\rho_1,
\qquad
t_M m_M=t_M\rho_2,
\]
where \(\rho_1,\rho_2\) are the two projections from
\[
M_1\times_{t_M,P,s_M}M_1.
\]
The monoid axioms follow from the identity and associativity laws in
\(\mathbb A\), as in the construction of \(\mathsf{Prog}_\Phi\).
Hence the recovered internal category is \(\mathsf{Prog}_\Phi\).
\end{proof}

\begin{remark}
\label{rem:coalgebra-does-not-replace-program-category}
The polynomial coalgebra does not replace the program category. It is a
curried total-response presentation of the same Mealy data. The program
category is recovered by flattening the coalgebra to its input-indexed
execution span.
\end{remark}

\section{Four loci of guardedness}
\label{sec:four-loci-of-guardedness}

The polynomial presentation is useful because it separates the different
places where guarded conditions may be imposed.

For a response profile
\[
c(x)(a)=(\beta_a,x_a),
\]
there are four natural loci of guardedness:
\[
\theta\hookrightarrow P
\qquad
\text{current-state guard},
\]
\[
R\hookrightarrow A_1
\qquad
\text{input-query guard},
\]
\[
O\hookrightarrow B_1
\qquad
\text{output-response guard},
\]
and
\[
\varphi\hookrightarrow P
\qquad
\text{continuation-state predicate}.
\]

The existential guarded one-step condition has the external form
\[
\theta(x)
\quad\text{and}\quad
\exists a:u\to p(x):
R(a)\wedge O(\beta_a)\wedge\varphi(x_a).
\]
The span semantics expresses this as a restricted execution diamond.
The polynomial semantics expresses it as a guarded predicate lifting on
the total response object, pulled back along the coalgebra.

This is the main bridge between guarded dynamic logic and the polynomial
calculus in this chapter.

\section{Guarded primitive programs}
\label{sec:guarded-primitive-programs}

We now return to predicate semantics. The guarded program operations in
this section are inherited from the span semantics of the previous
chapter, but the polynomial presentation gives a curried account of their
one-step predicate liftings.

Assume in this section that \(\mathcal E\) is regular when diamonds are
used.

Let
\[
\theta:\Theta\hookrightarrow P
\]
be a state guard,
\[
R\hookrightarrow A_1
\]
an input-arrow guard, and
\[
O\hookrightarrow B_1
\]
an output-arrow guard.

The primitive execution object is
\[
M_1:=A_1\times_{t_A,A_0,p}P.
\]
Write
\[
\pi_1:M_1\to A_1,
\qquad
\pi_2:M_1\to P
\]
for its projections.

Pull the three guards back to \(M_1\):
\[
\theta^\sharp:=\pi_2^\ast\theta,
\qquad
R^\sharp:=\pi_1^\ast R,
\qquad
O^\sharp:=\omega^\ast O.
\]
These pullbacks are displayed by
\[
\begin{tikzcd}[column sep=large,row sep=large]
\theta^\sharp
  \arrow[r, hook]
  \arrow[d]
&
M_1
  \arrow[d, "\pi_2"]
\\
\Theta
  \arrow[r, hook, "\theta"']
&
P
\end{tikzcd}
\qquad
\begin{tikzcd}[column sep=large,row sep=large]
R^\sharp
  \arrow[r, hook]
  \arrow[d]
&
M_1
  \arrow[d, "\pi_1"]
\\
R
  \arrow[r, hook]
&
A_1
\end{tikzcd}
\qquad
\begin{tikzcd}[column sep=large,row sep=large]
O^\sharp
  \arrow[r, hook]
  \arrow[d]
&
M_1
  \arrow[d, "\omega"]
\\
O
  \arrow[r, hook]
&
B_1 .
\end{tikzcd}
\]
Their simultaneous restriction is
\[
(\theta;R/O)^\sharp
:=
\theta^\sharp\wedge R^\sharp\wedge O^\sharp
\hookrightarrow M_1.
\]
Let
\[
i_{\theta;R/O}:(\theta;R/O)^\sharp\hookrightarrow M_1
\]
be the inclusion.

\begin{definition}[Guarded primitive Mealy program]
\label{def:guarded-primitive-mealy-program}
The \textbf{guarded primitive program}
\[
\theta;R/O
\]
is the restricted execution span
\[
\begin{tikzcd}[column sep=large]
P
&
(\theta;R/O)^\sharp
  \arrow[l, "\pi_2 i_{\theta;R/O}"']
  \arrow[r, "\delta i_{\theta;R/O}"]
&
P .
\end{tikzcd}
\]
Its diamond modality is denoted
\[
\langle \theta;R/O\rangle\varphi.
\]
\end{definition}

\begin{proposition}[External reading of guarded primitive programs]
\label{prop:external-reading-guarded-primitive}
In \(\mathbf{Set}\),
\[
x\models \langle \theta;R/O\rangle\varphi
\]
holds precisely when
\[
\theta(x)
\]
and there exists an admissible input arrow
\[
a:u\to p(x)
\]
such that
\[
R(a),
\qquad
O(\omega(a,x)),
\qquad
\delta(a,x)\models \varphi.
\]
\end{proposition}

\begin{proof}
By definition,
\[
\langle \theta;R/O\rangle\varphi
\]
is the diamond for the restricted span whose apex is
\[
(\theta;R/O)^\sharp.
\]
Thus \(x\) satisfies it precisely when there is a primitive execution
\[
(a,x)\in M_1
\]
such that
\[
\theta^\sharp(a,x),
\qquad
R^\sharp(a,x),
\qquad
O^\sharp(a,x),
\]
and whose target state
\[
\delta(a,x)
\]
satisfies \(\varphi\). These conditions say respectively
\[
\theta(x),
\qquad
R(a),
\qquad
O(\omega(a,x)),
\qquad
\varphi(\delta(a,x)).
\]
\end{proof}

\section{Guarded polynomial response liftings}
\label{sec:guarded-polynomial-response-liftings}

Let
\[
c:P\to\mathcal M_{\mathbb A,\mathbb B}(P)
\]
be the multiplicative Mealy coalgebra associated to
\[
(P,\delta,\omega):\mathbb A\rightsquigarrow\mathbb B.
\]

The coalgebraic form allows us to define guarded predicate liftings on the
total response object
\[
\mathcal M_{\mathbb A,\mathbb B}(P).
\]
This is in the spirit of coalgebraic modal logic, where modal operators
are interpreted by predicate liftings for the underlying coalgebraic
functor
\cite{PattinsonCoalgebraicModalLogic}.
Care is needed: a state guard
\[
\theta\hookrightarrow P
\]
is not a predicate on
\[
\mathcal M_{\mathbb A,\mathbb B}(P).
\]
Thus state guards are applied after pulling back along the coalgebra,
while input and output guards can be included directly in the response
lifting.

Let
\[
H_P:=
\Pi_{p_{\mathbb A,\mathbb B}}
s_{\mathbb A,\mathbb B}^{\ast}P.
\]
The object \(H_P\) is an object over \(K_{\mathbb A,\mathbb B}\). The
object
\[
\mathcal M_{\mathbb A,\mathbb B}(P)
=
\Sigma_{t_{\mathbb A,\mathbb B}}H_P
\]
has the same underlying object as \(H_P\), but regarded over \(S\) by
postcomposition with
\[
t_{\mathbb A,\mathbb B}:K_{\mathbb A,\mathbb B}\to S.
\]
We freely identify their underlying objects when defining the position
object.

Define the object of positions in a response profile by
\[
\mathrm{Pos}_P
:=
E_{\mathbb A,\mathbb B}\times_{K_{\mathbb A,\mathbb B}}H_P.
\]
Thus
\[
\begin{tikzcd}[column sep=large,row sep=large]
\mathrm{Pos}_P
  \arrow[r, "\rho"]
  \arrow[d]
&
H_P
  \arrow[d]
\\
E_{\mathbb A,\mathbb B}
  \arrow[r, "p_{\mathbb A,\mathbb B}"']
&
K_{\mathbb A,\mathbb B}
\end{tikzcd}
\]
is a pullback. The projection
\[
\rho:\mathrm{Pos}_P\to H_P\cong\mathcal M_{\mathbb A,\mathbb B}(P)
\]
forgets the chosen input position and remembers only the total response
profile.

There are canonical evaluation maps
\[
\operatorname{in}:\mathrm{Pos}_P\to A_1,
\]
\[
\operatorname{out}:\mathrm{Pos}_P\to B_1,
\]
and
\[
\operatorname{next}:\mathrm{Pos}_P\to P.
\]
They are induced as follows:
\[
\begin{tikzcd}[column sep=large]
\mathrm{Pos}_P
  \arrow[r]
  \arrow[rrr, bend left=15, "{\operatorname{in}}"]
&
E_{\mathbb A,\mathbb B}
  \arrow[r, "{\epsilon_Q}"]
&
Q_{\mathbb A}
  \arrow[r, "{\overline a}"]
&
A_1 .
\end{tikzcd}
\]
\[
\begin{tikzcd}[column sep=large]
\mathrm{Pos}_P
  \arrow[r]
  \arrow[rrr, bend left=15, "{\operatorname{out}}"]
&
E_{\mathbb A,\mathbb B}
  \arrow[r, "\operatorname{ev}"]
&
O_{\mathbb A,\mathbb B}=A_1\times B_1
  \arrow[r, "\pi_{B_1}"]
&
B_1 .
\end{tikzcd}
\]
The map
\[
\operatorname{next}:\mathrm{Pos}_P\to P
\]
is the counit evaluation of
\[
p_{\mathbb A,\mathbb B}^{\ast}H_P
\to
s_{\mathbb A,\mathbb B}^{\ast}P.
\]
It lies over the continuation-interface map:
\[
\begin{tikzcd}[column sep=large,row sep=large]
\mathrm{Pos}_P
  \arrow[r, "\operatorname{next}"]
  \arrow[d]
&
P
  \arrow[d, "{(p,q)}"]
\\
E_{\mathbb A,\mathbb B}
  \arrow[r, "s_{\mathbb A,\mathbb B}"']
&
S .
\end{tikzcd}
\]

\begin{definition}[Guarded existential response lifting]
\label{def:guarded-existential-response-lifting}
Assume \(\mathcal E\) is regular. Let
\[
R\hookrightarrow A_1,
\qquad
O\hookrightarrow B_1,
\qquad
\varphi\hookrightarrow P
\]
be predicates. Define
\[
\lambda_{R/O}(\varphi)
\hookrightarrow
\mathcal M_{\mathbb A,\mathbb B}(P)
\]
by
\[
\lambda_{R/O}(\varphi)
:=
\exists_{\rho}
\left(
\operatorname{in}^{\ast}R
\wedge
\operatorname{out}^{\ast}O
\wedge
\operatorname{next}^{\ast}\varphi
\right).
\]
Externally, for
\[
z\in \mathcal M_{\mathbb A,\mathbb B}(P)_{(v,y)},
\]
this says that
\[
z\in \lambda_{R/O}(\varphi)
\]
if and only if there exists an input arrow
\[
a:u\to v
\]
such that \(R(a)\), and, writing
\[
z(a)=(\beta,x'),
\]
we have
\[
O(\beta),
\qquad
\varphi(x').
\]
\end{definition}

\begin{proposition}[Agreement of guarded polynomial and span modalities]
\label{prop:guarded-polynomial-span-agreement}
Assume \(\mathcal E\) is regular, so that regular images satisfy
Beck--Chevalley and Frobenius reciprocity. Let
\[
c:P\to\mathcal M_{\mathbb A,\mathbb B}(P)
\]
be the multiplicative Mealy coalgebra associated to
\[
(P,\delta,\omega):\mathbb A\rightsquigarrow\mathbb B.
\]
Then, for every state guard
\[
\theta\hookrightarrow P,
\]
input guard
\[
R\hookrightarrow A_1,
\]
output guard
\[
O\hookrightarrow B_1,
\]
and postcondition
\[
\varphi\hookrightarrow P,
\]
we have
\[
\theta\wedge c^\ast\lambda_{R/O}(\varphi)
=
\langle \theta;R/O\rangle\varphi.
\]
\end{proposition}

\begin{proof}
Pull back the position object
\[
\rho:\mathrm{Pos}_P\to \mathcal M_{\mathbb A,\mathbb B}(P)
\]
along the coalgebra
\[
c:P\to\mathcal M_{\mathbb A,\mathbb B}(P).
\]
The resulting object is canonically the primitive execution object
\[
M_1=A_1\times_{t_A,A_0,p}P:
\]
\[
\begin{tikzcd}[column sep=large,row sep=large]
M_1
  \arrow[r]
  \arrow[d, "\pi_2"']
&
\mathrm{Pos}_P
  \arrow[d, "\rho"]
\\
P
  \arrow[r, "c"']
&
\mathcal M_{\mathbb A,\mathbb B}(P).
\end{tikzcd}
\]
Under this identification, the input, output, and continuation evaluations
are respectively
\[
\pi_1:M_1\to A_1,
\qquad
\omega:M_1\to B_1,
\qquad
\delta:M_1\to P.
\]
Therefore, by Beck--Chevalley for regular images,
\[
c^\ast\lambda_{R/O}(\varphi)
=
\exists_{\pi_2}
\left(
\pi_1^\ast R
\wedge
\omega^\ast O
\wedge
\delta^\ast\varphi
\right).
\]
Using Frobenius reciprocity for regular images,
\[
\theta\wedge c^\ast\lambda_{R/O}(\varphi)
=
\exists_{\pi_2}
\left(
\pi_2^\ast\theta
\wedge
\pi_1^\ast R
\wedge
\omega^\ast O
\wedge
\delta^\ast\varphi
\right).
\]
The right-hand side is precisely the diamond of the restricted execution
span
\[
\theta;R/O.
\]
Therefore
\[
\theta\wedge c^\ast\lambda_{R/O}(\varphi)
=
\langle \theta;R/O\rangle\varphi.
\]
\end{proof}

\begin{remark}[The central comparison]
\label{rem:central-guarded-comparison}
This comparison is the conceptual hinge of the chapter. The polynomial
coalgebra gives a guarded one-step response lifting on total response
profiles. Pulling that lifting back along the Mealy coalgebra and applying
the state guard recovers the ordinary guarded span diamond.
\end{remark}

\section{Profile modalities}
\label{sec:profile-modalities}

The existential lifting of
Definition~\ref{def:guarded-existential-response-lifting} recovers the
ordinary guarded span diamond after pulling back along the coalgebra. The
polynomial presentation also displays modalities that use the total
response profile more directly.

Let
\[
c(x)(a)=(\beta_a,x_a).
\]
For an input guard
\[
R\hookrightarrow A_1
\]
and output guard
\[
O\hookrightarrow B_1,
\]
one may define a universal profile condition
\[
\forall a\in R:
\quad
O(\beta_a)\Rightarrow \varphi(x_a).
\]
This treats \(O\) as an output filter: whenever the produced output lies
in \(O\), the corresponding successor must satisfy \(\varphi\). Internally,
in a first-order Heyting setting, this is represented by
\[
\Box^{\mathrm{prof}}_{R/O}(\varphi)
:=
\forall_\rho
\left(
(\operatorname{in}^{\ast}R\wedge \operatorname{out}^{\ast}O)
\Rightarrow
\operatorname{next}^{\ast}\varphi
\right).
\]

A stronger specification-style condition is
\[
\forall a\in R:
\quad
O(\beta_a)\wedge \varphi(x_a).
\]
This requires every guarded input to produce an \(O\)-output and land in a
\(\varphi\)-state. Internally, one may write
\[
\Box^{\mathrm{prof,strong}}_{R/O}(\varphi)
:=
\forall_\rho
\left(
\operatorname{in}^{\ast}R
\Rightarrow
\bigl(
\operatorname{out}^{\ast}O
\wedge
\operatorname{next}^{\ast}\varphi
\bigr)
\right).
\]

The distinction is visible in the contrast between the coalgebraic
profile diagram
\[
\begin{tikzcd}[column sep=large]
\mathrm{Pos}_P
  \arrow[r, "\rho"]
&
\mathcal M_{\mathbb A,\mathbb B}(P)
&
P
  \arrow[l, "c"']
\end{tikzcd}
\]
and the flattened execution span
\[
\begin{tikzcd}[column sep=large]
P
&
M_1
  \arrow[l, "\pi_2"']
  \arrow[r, "\delta"]
&
P .
\end{tikzcd}
\]
Profile modalities quantify over positions of the total response profile,
whereas span modalities quantify over execution witnesses.

\begin{remark}[Profile modalities versus span boxes]
\label{rem:profile-modalities-versus-span-boxes}
Profile modalities quantify over input positions in the total response
profile
\[
c(x):
\prod_{a:u\to p(x)}
\sum_{\beta:y'\to q(x)}
P_{(u,y')}.
\]
This is different from the box modality of a restricted execution span.
The span box quantifies over execution witnesses belonging to a restricted
span; a profile modality quantifies over input positions of a total
response policy. These coincide only under additional assumptions about
how the restriction is interpreted.
\end{remark}

\section{Tests, guarded commands, and guarded choice}
\label{sec:tests-guarded-commands-guarded-choice}

The remaining guarded program operations are inherited from the span
dynamic logic of Chapter~\ref{ch:mealy-program-categories-dynamic-logic}.
The polynomial construction explains the one-step guarded response
liftings; tests, conditionals, choice, and while programs are then formed
in the span calculus. The terminology of guarded commands follows
Dijkstra's guarded-command calculus \cite{DijkstraGuardedCommands}.

Assume now that \(\mathcal E\) is regular. Let
\[
\theta:\Theta\hookrightarrow P
\]
be a state predicate. The associated test span is
\[
\theta?
=
\left(
P\xleftarrow{\theta}\Theta\xrightarrow{\theta}P
\right),
\]
displayed as
\[
\begin{tikzcd}[column sep=large]
P
&
\Theta
  \arrow[l, "\theta"']
  \arrow[r, "\theta"]
&
P .
\end{tikzcd}
\]

If
\[
T=
\left(
P\xleftarrow{\ell_T}T\xrightarrow{r_T}P
\right)
\]
is a fixed ambient endoprogram and
\[
i_U:U\hookrightarrow T
\]
is a subprogram, then the guarded command
\[
\theta?;U
\]
is the composite span
\[
P\xleftarrow{} \Theta\times_{P}U\xrightarrow{}P,
\]
where the pullback is taken over
\[
\theta:\Theta\to P
\qquad
\text{and}
\qquad
\ell_T i_U:U\to P.
\]
Thus
\[
\begin{tikzcd}[column sep=large,row sep=large]
\Theta\times_P U
  \arrow[r]
  \arrow[d]
&
U
  \arrow[d, "\ell_T i_U"]
\\
\Theta
  \arrow[r, "\theta"']
&
P
\end{tikzcd}
\]
is a pullback. Under the inclusion \(U\hookrightarrow T\), this composite
is canonically identified with the subobject
\[
\ell_T^\ast\theta\wedge U
\hookrightarrow T.
\]

\begin{proposition}[Guard-prefix law]
\label{prop:guard-prefix-law}
In a regular category, for every subprogram
\[
U\hookrightarrow T
\]
of a fixed ambient endospan \(T:P\nrightarrow P\),
\[
\langle \theta?;U\rangle\varphi
=
\theta\wedge \langle U\rangle\varphi.
\]
\end{proposition}

\begin{proof}
By functoriality of span diamonds,
\[
\langle \theta?;U\rangle\varphi
=
\langle \theta?\rangle(\langle U\rangle\varphi).
\]
The test law gives
\[
\langle \theta?\rangle(\langle U\rangle\varphi)
=
\theta\wedge\langle U\rangle\varphi.
\]
\end{proof}

\begin{remark}[Guarded commands as test-prefixing]
\label{rem:guarded-commands-as-test-prefixing}
Guarding is test-prefixing:
\[
\theta\triangleright U
\quad
\text{means}
\quad
\theta?;U.
\]
For subprograms of a fixed ambient span
\[
T=
\left(P\xleftarrow{\ell_T}T\xrightarrow{r_T}P\right),
\]
this composite is canonically identified with the restriction by the
source-state predicate:
\[
\theta?;U
\cong
(\ell_T^\ast\theta)\wedge U
\hookrightarrow T.
\]
In particular, for primitive subprograms
\[
U\hookrightarrow M_1,
\]
where \(\ell_T=\pi_2\), this becomes
\[
\theta?;U
\cong
(\pi_2^\ast\theta)\wedge U
\hookrightarrow M_1.
\]
\end{remark}

Assume now that \(\mathcal E\) is coherent, so that finite joins of
subobjects exist and are stable under pullback.

Let
\[
T:P\nrightarrow P
\]
be a fixed ambient endospan, let
\[
U,V\hookrightarrow T
\]
be subprograms, and let
\[
\theta,\epsilon\hookrightarrow P
\]
be state guards. We regard
\[
\theta?;U
\qquad
\text{and}
\qquad
\epsilon?;V
\]
as subobjects of \(T\) via the identifications
\[
\theta?;U
\cong
(\ell_T^\ast\theta)\wedge U,
\]
\[
\epsilon?;V
\cong
(\ell_T^\ast\epsilon)\wedge V.
\]
Their guarded choice is the subprogram
\[
(\theta?;U)\vee(\epsilon?;V)
\hookrightarrow T.
\]
The common-ambient requirement is essential: coherent finite joins provide
joins of subobjects of a fixed object. Choice between arbitrary spans
requires additional coproduct structure and is a different, witnessful
operation.

\begin{proposition}[Diamond law for guarded choice]
\label{prop:diamond-law-guarded-choice}
For guarded branches
\[
\theta?;U
\qquad
\text{and}
\qquad
\epsilon?;V
\]
inside a fixed ambient endospan \(T:P\nrightarrow P\), we have
\[
\left\langle
(\theta?;U)\vee(\epsilon?;V)
\right\rangle\varphi
=
(\theta\wedge\langle U\rangle\varphi)
\vee
(\epsilon\wedge\langle V\rangle\varphi).
\]
\end{proposition}

\begin{proof}
By the finite choice law for subprograms in a common ambient span,
\[
\left\langle
(\theta?;U)\vee(\epsilon?;V)
\right\rangle\varphi
=
\langle \theta?;U\rangle\varphi
\vee
\langle \epsilon?;V\rangle\varphi.
\]
By Proposition~\ref{prop:guard-prefix-law},
\[
\langle \theta?;U\rangle\varphi
=
\theta\wedge\langle U\rangle\varphi
\]
and
\[
\langle \epsilon?;V\rangle\varphi
=
\epsilon\wedge\langle V\rangle\varphi.
\]
Combining these gives the displayed equality.
\end{proof}

If the guards satisfy
\[
\theta\wedge\epsilon=\bot,
\qquad
\theta\vee\epsilon=\top,
\]
then the guarded choice behaves as a deterministic branch at the level of
guard coverage. In Boolean predicate fibres one may take
\[
\epsilon=\neg\theta.
\]
In an intuitionistic setting, it is preferable to work with an explicitly
chosen complementary exit guard \(\epsilon\).

\begin{definition}[Internal conditional]
\label{def:internal-conditional}
Let
\[
T:P\nrightarrow P
\]
be a fixed ambient endospan, and let
\[
U,V\hookrightarrow T
\]
be subprograms. Let
\[
\theta,\epsilon\hookrightarrow P
\]
be complementary guards. Define
\[
\mathbf{if}\ \theta\ \mathbf{then}\ U\ \mathbf{else}\ V
:=
(\theta?;U)\vee(\epsilon?;V)
\hookrightarrow T.
\]
In Boolean fibres this may be written
\[
(\theta?;U)\vee((\neg\theta)?;V).
\]
\end{definition}

\begin{corollary}[Diamond law for conditionals]
\label{cor:conditional-diamond-law}
For complementary guards \(\theta,\epsilon\) and subprograms
\[
U,V\hookrightarrow T
\]
of a fixed ambient endospan,
\[
\left\langle
\mathbf{if}\ \theta\ \mathbf{then}\ U\ \mathbf{else}\ V
\right\rangle\varphi
=
(\theta\wedge\langle U\rangle\varphi)
\vee
(\epsilon\wedge\langle V\rangle\varphi).
\]
\end{corollary}

\begin{proof}
This is Proposition~\ref{prop:diamond-law-guarded-choice}.
\end{proof}

\section{Boxes for guarded choice and conditionals}
\label{sec:boxes-guarded-choice-conditionals}

Assume now that \(\mathcal E\) is a first-order Heyting category.

Let
\[
T=
\left(
P\xleftarrow{\ell_T}T\xrightarrow{r_T}P
\right)
\]
be a fixed ambient endospan. For a subprogram
\[
W\hookrightarrow T
\]
write
\[
[W]\varphi
\]
for the box modality of the restricted span determined by \(W\).

The box modality for a test is
\[
[\theta?]\varphi=\theta\Rightarrow\varphi.
\]
Therefore guarded prefixing gives
\[
[\theta?;U]\varphi
=
[\theta?]([U]\varphi)
=
\theta\Rightarrow[U]\varphi.
\]

\begin{proposition}[Box of finite subprogram choice]
\label{prop:box-finite-subprogram-choice}
For subprograms
\[
W,Z\hookrightarrow T
\]
of a fixed ambient endospan,
\[
[W\vee Z]\varphi
=
[W]\varphi\wedge[Z]\varphi.
\]
\end{proposition}

\begin{proof}
Regard \(W\) and \(Z\) as predicates on the apex of \(T\). Then
\[
[W]\varphi
=
\forall_{\ell_T}\bigl(W\Rightarrow r_T^\ast\varphi\bigr).
\]
Indeed, if \(i_W:W\hookrightarrow T\), then
\[
[W]\varphi
=
\forall_{\ell_T i_W}(i_W^\ast r_T^\ast\varphi)
=
\forall_{\ell_T}\forall_{i_W}(i_W^\ast r_T^\ast\varphi),
\]
and for a monomorphism \(i_W\),
\[
\forall_{i_W}(i_W^\ast A)=W\Rightarrow A.
\]
Therefore
\[
[W\vee Z]\varphi
=
\forall_{\ell_T}\bigl((W\vee Z)\Rightarrow r_T^\ast\varphi\bigr).
\]
In a Heyting algebra,
\[
(W\vee Z)\Rightarrow A
=
(W\Rightarrow A)\wedge(Z\Rightarrow A).
\]
Since \(\forall_{\ell_T}\) is a right adjoint, it preserves finite meets.
Hence
\[
[W\vee Z]\varphi
=
[W]\varphi\wedge[Z]\varphi.
\]
\end{proof}

\begin{proposition}[Box law for guarded conditionals]
\label{prop:box-law-guarded-conditionals}
For complementary guards \(\theta,\epsilon\) and subprograms
\[
U,V\hookrightarrow T
\]
of a fixed ambient endospan,
\[
\left[
\mathbf{if}\ \theta\ \mathbf{then}\ U\ \mathbf{else}\ V
\right]\varphi
=
(\theta\Rightarrow[U]\varphi)
\wedge
(\epsilon\Rightarrow[V]\varphi).
\]
\end{proposition}

\begin{proof}
The conditional is
\[
(\theta?;U)\vee(\epsilon?;V).
\]
By Proposition~\ref{prop:box-finite-subprogram-choice},
\[
\left[
(\theta?;U)\vee(\epsilon?;V)
\right]\varphi
=
[\theta?;U]\varphi
\wedge
[\epsilon?;V]\varphi.
\]
Finally,
\[
[\theta?;U]\varphi
=
\theta\Rightarrow[U]\varphi
\]
and similarly for \(\epsilon?;V\).
\end{proof}

\section{Guarded iteration and while programs}
\label{sec:guarded-iteration-while-programs}

Assume now that \(\mathcal E\) has the path-star structure used in
Chapter~\ref{ch:mealy-program-categories-dynamic-logic}. Let
\[
U:P\nrightarrow P
\]
be an endoprogram and let
\[
\theta,\epsilon\hookrightarrow P
\]
be guards. The intended reading is that \(\theta\) is the loop guard and
\(\epsilon\) is the exit guard.

Define the guarded body
\[
\theta?;U.
\]
The while program with exit guard \(\epsilon\) is
\[
W_{\theta,U,\epsilon}
:=
(\theta?;U)^\star_{\mathrm{path}};\epsilon?.
\]

\begin{definition}[Guarded while program]
\label{def:guarded-while-program}
The guarded while program with loop guard \(\theta\), body \(U\), and exit
guard \(\epsilon\) is
\[
\mathbf{while}\ \theta\ \mathbf{do}\ U\ \mathbf{exit}\ \epsilon
:=
(\theta?;U)^\star_{\mathrm{path}};\epsilon?.
\]
In Boolean predicate fibres, the ordinary notation
\[
\mathbf{while}\ \theta\ \mathbf{do}\ U
\]
denotes the special case
\[
(\theta?;U)^\star_{\mathrm{path}};(\neg\theta)?.
\]
\end{definition}

The form of the while laws is the usual dynamic-logic and
Kleene-algebraic reading of iteration and tests, here interpreted in the
span semantics \cite{HarelKozenTiurynDynamicLogic,KozenKAT}.

\begin{theorem}[Diamond law for guarded while programs]
\label{thm:diamond-law-guarded-while}
In an \(\omega\)-iterative first-order Heyting category,
\[
\left\langle
(\theta?;U)^\star_{\mathrm{path}};\epsilon?
\right\rangle\varphi
=
\mu X.\bigl(\epsilon\wedge\varphi
\vee
\theta\wedge\langle U\rangle X\bigr).
\]
\end{theorem}

\begin{proof}
By functoriality of diamonds,
\[
\left\langle
(\theta?;U)^\star_{\mathrm{path}};\epsilon?
\right\rangle\varphi
=
\left\langle
(\theta?;U)^\star_{\mathrm{path}}
\right\rangle
(\langle \epsilon?\rangle\varphi).
\]
The test law gives
\[
\langle \epsilon?\rangle\varphi
=
\epsilon\wedge\varphi.
\]
By the path/fixed-point comparison theorem from
Chapter~\ref{ch:mealy-program-categories-dynamic-logic},
\[
\left\langle
(\theta?;U)^\star_{\mathrm{path}}
\right\rangle
(\epsilon\wedge\varphi)
=
\mu X.\bigl((\epsilon\wedge\varphi)
\vee
\langle \theta?;U\rangle X\bigr).
\]
By the guard-prefix law,
\[
\langle \theta?;U\rangle X
=
\theta\wedge\langle U\rangle X.
\]
Therefore
\[
\left\langle
(\theta?;U)^\star_{\mathrm{path}};\epsilon?
\right\rangle\varphi
=
\mu X.\bigl(\epsilon\wedge\varphi
\vee
\theta\wedge\langle U\rangle X\bigr).
\]
\end{proof}

\begin{corollary}[Boolean while law]
\label{cor:boolean-while-law}
In Boolean predicate fibres,
\[
\left\langle
\mathbf{while}\ \theta\ \mathbf{do}\ U
\right\rangle\varphi
=
\mu X.\bigl((\neg\theta\wedge\varphi)
\vee
(\theta\wedge\langle U\rangle X)\bigr).
\]
\end{corollary}

\begin{proof}
Take
\[
\epsilon=\neg\theta
\]
in Theorem~\ref{thm:diamond-law-guarded-while}.
\end{proof}

\begin{theorem}[Box law for guarded while programs]
\label{thm:box-law-guarded-while}
In an \(\omega\)-iterative first-order Heyting category,
\[
\left[
(\theta?;U)^\star_{\mathrm{path}};\epsilon?
\right]\varphi
=
\nu X.\bigl((\epsilon\Rightarrow\varphi)
\wedge
(\theta\Rightarrow[U]X)\bigr).
\]
\end{theorem}

\begin{proof}
By functoriality of boxes,
\[
\left[
(\theta?;U)^\star_{\mathrm{path}};\epsilon?
\right]\varphi
=
\left[
(\theta?;U)^\star_{\mathrm{path}}
\right]
([\epsilon?]\varphi).
\]
The box test law gives
\[
[\epsilon?]\varphi
=
\epsilon\Rightarrow\varphi.
\]
By the path/fixed-point comparison theorem for boxes,
\[
\left[
(\theta?;U)^\star_{\mathrm{path}}
\right]
(\epsilon\Rightarrow\varphi)
=
\nu X.\bigl((\epsilon\Rightarrow\varphi)
\wedge
[\theta?;U]X\bigr).
\]
Finally,
\[
[\theta?;U]X
=
\theta\Rightarrow[U]X.
\]
Thus
\[
\left[
(\theta?;U)^\star_{\mathrm{path}};\epsilon?
\right]\varphi
=
\nu X.\bigl((\epsilon\Rightarrow\varphi)
\wedge
(\theta\Rightarrow[U]X)\bigr).
\]
\end{proof}

\section{Domain and dynamic tests}
\label{sec:domain-and-dynamic-tests}

Let
\[
U:P\nrightarrow P
\]
be an endoprogram. Its domain predicate is
\[
\operatorname{dom}(U):=\langle U\rangle\top.
\]
This is the modal/domain operation familiar from Kleene algebra with
domain and antidomain calculi
\cite{DesharnaisMoellerStruthKleeneAlgebraDomain}.
Thus \(\operatorname{dom}(U)\) is the predicate of states from which at
least one \(U\)-execution is possible.

For a guarded primitive program
\[
\theta;R/O,
\]
the domain predicate is
\[
\operatorname{dom}(\theta;R/O)
=
\langle \theta;R/O\rangle\top.
\]
In generalized-element notation, this says that \(x\) lies in the domain
precisely when
\[
\theta(x)
\]
and there exists an admissible input arrow
\[
a:u\to p(x)
\]
such that
\[
R(a),
\qquad
O(\omega(a,x)).
\]

\begin{proposition}[Domain-prefix law]
\label{prop:domain-prefix-law}
For every endoprogram \(U\),
\[
\langle \operatorname{dom}(U)?;U\rangle\varphi
=
\langle U\rangle\varphi.
\]
\end{proposition}

\begin{proof}
By the guard-prefix law,
\[
\langle \operatorname{dom}(U)?;U\rangle\varphi
=
\operatorname{dom}(U)\wedge\langle U\rangle\varphi.
\]
Since
\[
\langle U\rangle\varphi\le \langle U\rangle\top
=
\operatorname{dom}(U),
\]
we have
\[
\operatorname{dom}(U)\wedge\langle U\rangle\varphi
=
\langle U\rangle\varphi.
\]
\end{proof}

If the predicate fibres are Boolean, define the antidomain predicate by
\[
\operatorname{adom}(U):=\neg\operatorname{dom}(U).
\]

\begin{proposition}[Antidomain annihilation]
\label{prop:antidomain-annihilation}
In Boolean predicate fibres,
\[
\langle \operatorname{adom}(U)?;U\rangle\varphi
=
\bot.
\]
\end{proposition}

\begin{proof}
By the guard-prefix law,
\[
\langle \operatorname{adom}(U)?;U\rangle\varphi
=
\operatorname{adom}(U)\wedge\langle U\rangle\varphi.
\]
Since
\[
\langle U\rangle\varphi\le \operatorname{dom}(U)
\]
and
\[
\operatorname{adom}(U)=\neg\operatorname{dom}(U),
\]
the meet is bottom.
\end{proof}

\section{Guarded Mealy automata}
\label{sec:guarded-mealy-automata}

We now isolate the finite guarded structure generated by a Mealy morphism.

Let
\[
(P,\delta,\omega):\mathbb A\rightsquigarrow\mathbb B
\]
be a Mealy morphism. Let
\[
(\theta_i)_{i\in I}
\]
be a finite family of state guards,
\[
(R_i)_{i\in I}
\]
a finite family of input-arrow guards, and
\[
(O_i)_{i\in I}
\]
a finite family of output-arrow guards. Each triple determines a guarded
primitive subprogram
\[
G_i:=(\theta_i;R_i/O_i)^\sharp\hookrightarrow M_1.
\]
These fit into the common ambient primitive span
\[
\begin{tikzcd}[column sep=large]
P
&
M_1
  \arrow[l, "\pi_2"']
  \arrow[r, "\delta"]
&
P ,
\end{tikzcd}
\]
so their finite join is formed in \(\mathrm{Sub}(M_1)\).

\begin{definition}[Guarded Mealy automaton]
\label{def:guarded-mealy-automaton}
A \textbf{guarded Mealy automaton} over
\[
(P,\delta,\omega)
\]
is a finite family of guarded primitive subprograms
\[
G_i=(\theta_i;R_i/O_i)^\sharp\hookrightarrow M_1.
\]
Its associated one-step program is the finite join
\[
G:=\bigvee_{i\in I}G_i\hookrightarrow M_1.
\]
\end{definition}

\begin{proposition}[Diamond semantics of guarded Mealy automata]
\label{prop:guarded-mealy-automaton-diamond}
For a guarded Mealy automaton with one-step program
\[
G=\bigvee_{i\in I}G_i,
\]
we have
\[
\langle G\rangle\varphi
=
\bigvee_{i\in I}
\langle G_i\rangle\varphi.
\]
Equivalently,
\[
\langle G\rangle\varphi
=
\bigvee_{i\in I}
\langle \theta_i;R_i/O_i\rangle\varphi.
\]
\end{proposition}

\begin{proof}
This is the finite choice law for subprograms applied to
\[
G=\bigvee_{i\in I}G_i
\]
inside the common ambient object \(M_1\).
\end{proof}

A guarded Mealy automaton is \textbf{witness-disjoint} if
\[
G_i\wedge G_j=\bot
\qquad
(i\ne j).
\]
It is \textbf{branch-disjoint} if its enabled-state predicates are
pairwise disjoint:
\[
\operatorname{dom}(G_i)\wedge \operatorname{dom}(G_j)=\bot
\qquad
(i\ne j).
\]
It is \textbf{guard-total} if
\[
\bigvee_{i\in I}\operatorname{dom}(G_i)=\top.
\]
When branch-disjointness and guard-totality hold, the enabled-state
predicates form an internally disjoint cover of \(P\). In \(\mathbf{Set}\),
this says that exactly one branch is enabled at each state. The enabled
branch may still contain several execution witnesses, so this is
deterministic only at the level of branch selection, not necessarily at
the level of primitive executions.

\section{Output-guarded loops and interface guards}
\label{sec:output-guarded-loops-interface-guards}

The Mealy setting allows guards to inspect emitted output arrows. This
gives guarded structure not present in ordinary state-only guarded command
languages.

Let
\[
U\hookrightarrow M_1
\]
be a primitive subprogram, and let
\[
O\hookrightarrow B_1
\]
be an output guard. Define the output-restricted body
\[
U_O:=U\wedge \omega^\ast O\hookrightarrow M_1.
\]
The pullback defining the output restriction is
\[
\begin{tikzcd}[column sep=large,row sep=large]
\omega^\ast O
  \arrow[r, hook]
  \arrow[d]
&
M_1
  \arrow[d, "\omega"]
\\
O
  \arrow[r, hook]
&
B_1 .
\end{tikzcd}
\]
Given loop and exit guards
\[
\theta,\epsilon\hookrightarrow P,
\]
define the output-guarded while program
\[
W_{\theta,U,O,\epsilon}
:=
(\theta?;U_O)^\star_{\mathrm{path}};\epsilon?.
\]

\begin{proposition}[Diamond law for output-guarded loops]
\label{prop:output-guarded-loop-diamond}
In an \(\omega\)-iterative first-order Heyting category,
\[
\langle W_{\theta,U,O,\epsilon}\rangle\varphi
=
\mu X.\bigl(
\epsilon\wedge\varphi
\vee
\theta\wedge\langle U_O\rangle X
\bigr).
\]
\end{proposition}

\begin{proof}
This is Theorem~\ref{thm:diamond-law-guarded-while} applied to the body
\[
U_O.
\]
\end{proof}

\begin{remark}[Output-guarded loops]
\label{rem:output-guarded-loop}
Output-guarded loops show why the Mealy setting is richer than an ordinary
guarded transition system. The continuation of the loop can depend not
only on the current state and input guard, but also on the class of output
arrows emitted by the primitive execution.
\end{remark}

Since every state has an input and output interface
\[
(p(x),q(x))\in A_0\times B_0,
\]
guards may also be placed directly on interface pairs.

Let
\[
\Gamma\hookrightarrow A_0\times B_0
\]
be an interface guard. Pulling it back along
\[
(p,q):P\to A_0\times B_0
\]
gives a state guard
\[
(p,q)^\ast\Gamma\hookrightarrow P:
\]
\[
\begin{tikzcd}[column sep=large,row sep=large]
(p,q)^\ast\Gamma
  \arrow[r, hook]
  \arrow[d]
&
P
  \arrow[d, "{(p,q)}"]
\\
\Gamma
  \arrow[r, hook]
&
A_0\times B_0 .
\end{tikzcd}
\]
Thus any guarded primitive program may be further restricted by interface
type:
\[
((p,q)^\ast\Gamma;R/O)^\sharp\hookrightarrow M_1.
\]

\begin{proposition}[External reading of interface-guarded primitives]
\label{prop:interface-guarded-external-reading}
In \(\mathbf{Set}\),
\[
x\models \langle (p,q)^\ast\Gamma;R/O\rangle\varphi
\]
holds precisely when
\[
\Gamma(p(x),q(x))
\]
and there exists an admissible input arrow
\[
a:u\to p(x)
\]
such that
\[
R(a),
\qquad
O(\omega(a,x)),
\qquad
\delta(a,x)\models\varphi.
\]
\end{proposition}

\begin{proof}
This is Proposition~\ref{prop:external-reading-guarded-primitive} with
\[
\theta=(p,q)^\ast\Gamma.
\]
\end{proof}

\begin{remark}[Typed guarded control]
\label{rem:typed-guarded-control}
Interface guards are one of the places where the internal typed setting
differs from ordinary untyped guarded command languages. The guard may
specify not only a region of the state object \(P\), but also the input
and output interfaces over which the state lies.
\end{remark}

\section{Guarded program algebra generated by a Mealy morphism}
\label{sec:guarded-program-algebra-generated-mealy}

The constructions above generate a guarded program calculus from a Mealy
morphism.

The primitive guarded commands are the guarded primitive programs
\[
\theta;R/O.
\]
They may be combined by:
\[
\text{sequential composition},
\qquad
\text{guarded finite choice inside a common ambient span},
\qquad
\text{tests},
\qquad
\text{guarded iteration}.
\]

The basic laws are:
\[
\langle \theta?;U\rangle\varphi
=
\theta\wedge\langle U\rangle\varphi,
\]
\[
\left\langle
(\theta?;U)\vee(\epsilon?;V)
\right\rangle\varphi
=
(\theta\wedge\langle U\rangle\varphi)
\vee
(\epsilon\wedge\langle V\rangle\varphi),
\]
where \(U\) and \(V\) are subprograms of a common ambient span, and
\[
\left\langle
(\theta?;U)^\star_{\mathrm{path}};\epsilon?
\right\rangle\varphi
=
\mu X.\bigl(
\epsilon\wedge\varphi
\vee
\theta\wedge\langle U\rangle X
\bigr).
\]

\begin{definition}[Internal guarded Mealy program algebra]
\label{def:internal-guarded-mealy-program-algebra}
The \textbf{internal guarded Mealy program algebra} generated by a Mealy
morphism
\[
(P,\delta,\omega):\mathbb A\rightsquigarrow\mathbb B
\]
is the class of endoprograms on \(P\) generated from:
\begin{enumerate}
\item tests \(\theta?\) for state predicates \(\theta\hookrightarrow P\);
\item guarded primitive programs \(\theta;R/O\);
\item sequential composition of spans;
\item finite guarded joins of subprograms inside a specified common
ambient span;
\item path-star guarded iteration
\[
(\theta?;U)^\star_{\mathrm{path}};\epsilon?.
\]
\end{enumerate}
\end{definition}

\begin{remark}[Common ambient spans versus arbitrary choice]
\label{rem:common-ambient-versus-arbitrary-choice}
Finite guarded joins in this definition are joins of subobjects inside a
specified ambient execution span. Thus the construction gives a guarded
subprogram calculus relative to chosen ambient spans. If one wants finite
choice between arbitrary generated spans, additional coproduct structure on
\(\mathcal E\) must be assumed and choice should be interpreted as
coproduct of span apexes. That is a different, witnessful choice
operation.
\end{remark}

\begin{remark}[Guarded program structure]
\label{rem:guarded-program-structure}
The guarded program algebra is obtained by transporting the span-based
dynamic logic of Chapter~\ref{ch:mealy-program-categories-dynamic-logic}
through the polynomial currying of the Mealy data. The genuinely
coalgebraic information is the total response profile
\[
c(x):
\prod_{a:u\to p(x)}
\sum_{\beta:y'\to q(x)}
P_{(u,y')},
\]
while the span program algebra is recovered by flattening this profile to
individual execution witnesses.
\end{remark}

\section{Representable carriers}
\label{sec:representable-carriers-polynomial-coalgebra}

Let
\[
F:\mathbb A\to\mathbb B
\]
be an internal functor with object map
\[
F_0:A_0\to B_0
\]
and arrow map
\[
F_1:A_1\to B_1.
\]
Regard \(F\) as the representable Mealy morphism carried by
\[
F_{0!}
=
\left(
A_0\xleftarrow{1_{A_0}}A_0\xrightarrow{F_0}B_0
\right).
\]
Thus
\[
P=A_0,
\qquad
p=1_{A_0},
\qquad
q=F_0.
\]
The carrier diagram is
\[
\begin{tikzcd}[column sep=large,row sep=large]
A_0
  \arrow[rr, "F_0"]
  \arrow[dr, "{(1_{A_0},F_0)}"']
  \arrow[ddr, "1_{A_0}"', bend right=12]
&&
B_0
\\
&
A_0\times B_0
  \arrow[ur, "\pi_B"']
  \arrow[d, "\pi_A"]
&
\\
&
A_0 .
\end{tikzcd}
\]

The associated coalgebra
\[
c_F:A_0\to\mathcal M_{\mathbb A,\mathbb B}(A_0)
\]
is given by
\[
c_F(v)(a)=(F_1(a),s_A(a))
\]
for every arrow
\[
a:u\to v
\]
of \(\mathbb A\).

\begin{theorem}[Representable Mealy coalgebras]
\label{thm:representable-mealy-coalgebras}
Let
\[
F:\mathbb A\to\mathbb B
\]
be an internal functor, regarded as a representable Mealy morphism. Then
its polynomial transpose is the coalgebra
\[
c_F:A_0\to\mathcal M_{\mathbb A,\mathbb B}(A_0)
\]
defined by
\[
c_F(v)(a)=(F_1(a),s_A(a)).
\]
Moreover, the program category recovered from \(c_F\) is canonically
isomorphic to
\[
\mathbb A^{\mathrm{op}},
\]
and the output labelling functor is
\[
F^{\mathrm{op}}:\mathbb A^{\mathrm{op}}\to\mathbb B^{\mathrm{op}}.
\]
\end{theorem}

\begin{proof}
For the representable carrier,
\[
P=A_0
\]
and a state is an object \(v\in A_0\). An admissible input arrow at \(v\)
is an arrow
\[
a:u\to v.
\]
The Mealy structure associated to the internal functor has
\[
\delta(a,v)=u=s_A(a)
\]
and
\[
\omega(a,v)=F_1(a).
\]
The typing of \(F_1(a)\) is given by the functoriality squares
\[
s_BF_1=F_0s_A,
\qquad
t_BF_1=F_0t_A.
\]
Thus
\[
F_1(a):F_0(u)\to F_0(v),
\]
so the polynomial transpose is
\[
c_F(v)(a)=(F_1(a),s_A(a)).
\]

By the result of Chapter~\ref{ch:mealy-program-categories-dynamic-logic},
the program category of a representable Mealy morphism is canonically
\[
\mathbb A^{\mathrm{op}},
\]
and the output labelling is
\[
F^{\mathrm{op}}.
\]
Since the coalgebra \(c_F\) recovers exactly this \(\delta\) and
\(\omega\), the same conclusion follows from the polynomial-coalgebraic
presentation.
\end{proof}

\section{Main consolidation}
\label{sec:polynomial-guarded-main-consolidation}

We collect the constructions of the chapter. The polynomial clauses below
require the products and dependent products used in the construction of
\(\mathcal M_{\mathbb A,\mathbb B}\). The diamond clauses require
regularity, the finite-choice clauses require coherence, and the box and
guarded-iteration clauses require the first-order Heyting and
\(\omega\)-iterative hypotheses specified in the relevant sections.

\begin{theorem}[Polynomial Mealy spans and guarded program structure]
\label{thm:polynomial-mealy-spans-guarded-consolidation}
Let
\[
\mathbb A
\qquad
\text{and}
\qquad
\mathbb B
\]
be internal categories in \(\mathcal E\). Under the corresponding
hypotheses just described, the following hold.

\begin{enumerate}
\item The pair \((\mathbb A,\mathbb B)\) determines a canonical Mealy
polynomial span over
\[
S=A_0\times B_0:
\]
\[
\begin{tikzcd}[column sep=large]
S
&
E_{\mathbb A,\mathbb B}
  \arrow[l, "s_{\mathbb A,\mathbb B}"']
  \arrow[r, "p_{\mathbb A,\mathbb B}"]
&
K_{\mathbb A,\mathbb B}
  \arrow[r, "t_{\mathbb A,\mathbb B}"]
&
S .
\end{tikzcd}
\]

\item This polynomial span induces an endofunctor
\[
\mathcal M_{\mathbb A,\mathbb B}:\mathcal E/S\to\mathcal E/S
\]
with fibre
\[
\mathcal M_{\mathbb A,\mathbb B}(X)_{(v,y)}
\cong
\prod_{a:u\to v}
\sum_{\beta:y'\to y}
X_{(u,y')}.
\]

\item For an object
\[
P\to S
\]
the typed one-step Mealy data
\[
\delta:
A_1\times_{t_A,A_0,p}P\to P,
\qquad
\omega:
A_1\times_{t_A,A_0,p}P\to B_1
\]
are equivalently coalgebra maps
\[
P\to\mathcal M_{\mathbb A,\mathbb B}(P).
\]

\item Under this equivalence, the Mealy unit and multiplication laws are
precisely the identity and multiplicativity laws of
Definition~\ref{def:multiplicative-mealy-coalgebra}. Hence Mealy
morphisms
\[
\mathbb A\rightsquigarrow\mathbb B
\]
are equivalently multiplicative Mealy coalgebras for
\[
\mathcal M_{\mathbb A,\mathbb B}.
\]

\item The primitive execution span and program category
\[
\mathsf{Prog}_\Phi
\]
of Chapter~\ref{ch:mealy-program-categories-dynamic-logic} are recovered
from the multiplicative coalgebra by flattening the total response profile
to the transition component
\[
\delta.
\]

\item State, input, and output guards
\[
\theta\hookrightarrow P,
\qquad
R\hookrightarrow A_1,
\qquad
O\hookrightarrow B_1
\]
determine guarded primitive programs
\[
\theta;R/O
\]
as subprograms of the primitive execution span.

\item The guarded span modalities agree with the guarded polynomial
response-lifting modalities after pulling back along the Mealy coalgebra:
\[
\theta\wedge c^\ast\lambda_{R/O}(\varphi)
=
\langle \theta;R/O\rangle\varphi.
\]

\item The polynomial presentation also supports profile modalities that
quantify over the input positions of the total response profile
\[
c(x):
\prod_{a:u\to p(x)}
\sum_{\beta:y'\to q(x)}
P_{(u,y')}.
\]
These profile modalities are not the same thing as boxes of restricted
execution spans, although the two notions may agree under additional
hypotheses.

\item The generated guarded program calculus satisfies the expected laws
\[
\langle \theta?;U\rangle\varphi
=
\theta\wedge\langle U\rangle\varphi,
\]
\[
\left\langle
(\theta?;U)\vee(\epsilon?;V)
\right\rangle\varphi
=
(\theta\wedge\langle U\rangle\varphi)
\vee
(\epsilon\wedge\langle V\rangle\varphi),
\]
where \(U\) and \(V\) are subprograms of a common ambient span, and
\[
\left\langle
(\theta?;U)^\star_{\mathrm{path}};\epsilon?
\right\rangle\varphi
=
\mu X.\bigl(
\epsilon\wedge\varphi
\vee
\theta\wedge\langle U\rangle X
\bigr).
\]

\item If the Mealy morphism is representable and induced by an internal
functor
\[
F:\mathbb A\to\mathbb B,
\]
then its polynomial transpose is
\[
c_F(v)(a)=(F_1(a),s_A(a)),
\]
and the recovered program category is
\[
\mathbb A^{\mathrm{op}}.
\]
\end{enumerate}
\end{theorem}

\begin{proof}
Part (1) is the construction of
Section~\ref{sec:queries-responses-canonical-mealy-polynomial}. Part (2)
is Proposition~\ref{prop:mealy-polynomial-fibre}. Part (3) is
Theorem~\ref{thm:polynomial-currying-one-step}. Part (4) is
Theorem~\ref{thm:mealy-morphisms-multiplicative-coalgebras}. Part (5) is
Theorem~\ref{thm:program-category-recovered-polynomial-coalgebra}. Part
(6) is Definition~\ref{def:guarded-primitive-mealy-program}. Part (7) is
Proposition~\ref{prop:guarded-polynomial-span-agreement}. Part (8) is
Section~\ref{sec:profile-modalities}. Part (9) follows from
Proposition~\ref{prop:guard-prefix-law},
Proposition~\ref{prop:diamond-law-guarded-choice}, and
Theorem~\ref{thm:diamond-law-guarded-while}. Part (10) is
Theorem~\ref{thm:representable-mealy-coalgebras}.
\end{proof}

\begin{remark}[Final summary]
\label{rem:final-polynomial-guarded-summary}
The polynomial structure in this chapter is not imposed from outside the
span calculus. It is obtained by currying the total one-step response
profile already present in a Mealy morphism:
\[
x
\longmapsto
\left(
a
\longmapsto
(\omega(a,x),\delta(a,x))
\right).
\]
The resulting coalgebra
\[
P\to\mathcal M_{\mathbb A,\mathbb B}(P)
\]
is the coalgebraic form of the same Mealy data that produced the execution
category in Chapter~\ref{ch:mealy-program-categories-dynamic-logic}.
Flattening the coalgebra recovers the primitive execution span and hence
the existing dynamic span semantics. The extra polynomial content is the
total response profile, which supports guarded response liftings and
profile modalities beyond ordinary span diamonds and boxes.
\end{remark}